\documentclass[11pt]{article}
\usepackage{fontspec}
\setmainfont{Times New Roman}
\setsansfont{Arial}
\setmonofont{Consolas}
\usepackage{amsmath,amssymb}
\usepackage{geometry}
\usepackage{microtype}
\geometry{a4paper,margin=3cm}
\setlength{\parindent}{1.25em}
\setlength{\parskip}{0pt}
\emergencystretch=30em
\allowdisplaybreaks
\raggedbottom
\providecommand{\Rfield}{\mathfrak R}
\providecommand{\bdelta}{\bar{\delta}}

\begin{document}

% Unit: Paper 12. Source segments P12-S0001--P12-S0019.

\begin{center}
{\Large\bfseries 12. Инварианты произвольных дифференциальных выражений}\par
\vspace{1.5em}
Nachr. v. d. Ges. d. Wiss. zu Göttingen 1918, с. 37--44
\end{center}

\vspace{3em}
\begin{center}
Представлено на заседании 25 января 1918 года Ф. Клейном.
\end{center}

Под дифференциальным выражением я понимаю функцию
\[
f(x,dx)=f(x_1,\ldots,x_n;\,dx_1,\ldots,dx_n)
\]
от \(n\) переменных \(x_1,\ldots,x_n\) и их дифференциалов
\(dx_1,\ldots,dx_n\), которая предполагается аналитической в обеих системах по \(n\) аргументов, но не обязательно вещественной. Теперь за основу я беру группу всех аналитических преобразований переменных, которой соответствует группа всех линейных преобразований дифференциалов:
\begin{equation}
x_i=x_i(y_1,\ldots,y_n);\qquad
dx_i=\sum_k\frac{\partial x_i}{\partial y_k}\,dy_k;\qquad
\delta x_i=\sum_k\frac{\partial x_i}{\partial y_k}\,\delta y_k,
\tag{1}
\end{equation}
при этом \(f(x,dx)\) пусть переходит в \(g(y,dy)\). Под инвариантом \(f(x,dx)\) я понимаю инвариант относительно этой группы и относительно индуцированной ею группы для высших дифференциалов \(d^2x,ddx,\ldots\) и для производных от \(f(x,dx)\); то есть такую аналитическую функцию \(J\), что для произвольного \(f\) и в силу (1) тождественно по переменным и всем входящим дифференциалам выполняется
\[
\begin{aligned}
&J\!\left(f,\frac{\partial f}{\partial dx}\cdots
\frac{\partial^{\rho+\sigma}f}{\partial x^\rho\partial dx^\sigma}
\cdots dx,\delta x,d^2x,\ldots\right)\\
&\qquad =
J\!\left(g,\frac{\partial g}{\partial dy}\cdots
\frac{\partial^{\rho+\sigma}g}{\partial y^\rho\partial dy^\sigma}
\cdots dy,\delta y,d^2y,\ldots\right).\footnotemark
\end{aligned}
\]
\footnotetext{\(\displaystyle
\frac{\partial^{\rho+\sigma}f}{\partial x^\rho\partial dx^\sigma}\)
всюду употребляется сокращенно для
\[
\frac{\partial^{\rho+\sigma}f}
{\partial x_{i_1}\cdots\partial x_{i_\rho}
 \partial dx_{k_1}\cdots\partial dx_{k_\sigma}},
\]
где индексы означают какие-либо из чисел от \(1\) до \(n\).}
Совершенно соответственно определяется инвариант одновременной системы дифференциальных выражений. Если такой инвариант, в частности, содержит только первые дифференциалы \(dx,\delta x\) и только производные по этим дифференциалам, а не по переменным, то появление переменных в функциях и линейное преобразование дифференциалов не принимаются во внимание; эти специальные инварианты являются инвариантами относительно группы всех линейных преобразований дифференциалов \(dx,\delta x\) с неопределенными коэффициентами и далее будут называться проективными инвариантами одновременной системы.

Для квадратичных однородных дифференциальных форм Кристоффель (Crelle 70) и Риччи (Math. Annalen 54) построили такую систему инвариантных образований, что все инварианты становятся проективными инвариантами этой одновременной системы; тем самым вопросы о всей совокупности инвариантов и об эквивалентности сводятся к вопросам линейной теории инвариантов. Это я хотела бы назвать «теоремой редукции». Полная система состоит из счетно бесконечного множества форм; но для инвариантов каждого конечного порядка \(\rho\), то есть таких, которые не содержат производных по переменным выше \(\rho\)-го порядка, она содержит только конечное число форм.

Ниже я показываю справедливость «теоремы редукции» для произвольных дифференциальных выражений; и здесь снова речь идет о полной системе счетно бесконечного множества инвариантных образований, но для каждого конечного порядка только о конечном их числе. Однако, тогда как Кристоффель и Риччи для порождения инвариантов и доказательства теоремы редукции опираются на трудно обозримые элиминации, я порождаю инварианты непосредственно инвариантными процессами дифференцирования и варьирования; последние просто можно понимать как процессы дифференцирования по другим параметрам. Алгоритм этих инвариантных процессов дифференцирования, в связи с Лагранжем (Mécanique analytique, 2-я часть, разд. IV), развили Риман (Werke XXII) и Липшиц (Crelle 70, 72) для квадратичных дифференциальных форм, однако не дойдя до теоремы редукции. Вспомогательное средство для доказательства теоремы редукции дают введенные Риманом для квадратичных форм «нормальные координаты» (Werke XIII), которые преобразуют экстремали соответствующей вариационной задачи, исходящие из одной точки, в прямые; но они существуют только для однородных дифференциальных выражений
\(f(x,\varkappa\cdot dx)=\Phi(\varkappa)\cdot f(x,dx)\footnote{Легко видеть, что \(\Phi(\varkappa)=\varkappa^c\), где \(c\) обозначает произвольную вещественную или комплексную величину.}\). Неоднородный случай, однако, сводится к этому.

\begin{center}
\textbf{I. Порождение инвариантов.}
\end{center}

Последовательность проективных инвариантов \(f(x,dx)\), вообще говоря не обрывающаяся, дается полярами. Из-за линейности преобразования дифференциалов из \(f(x,dx)=g(y,dy)\) также следует
\(f(x,dx+\lambda\delta x)=g(y,dy+\lambda\delta y)\); а отсюда, разлагая по \(\lambda\), получаются характерные для инвариантности соотношения:
\begin{equation}
\begin{gathered}
f(dx)=g(dy),\\
f_\delta=\sum\frac{\partial f(dx)}{\partial dx}\,\delta x
       =\sum\frac{\partial g(dy)}{\partial dy}\,\delta y,\\
\vdots\\
f_{\delta^\sigma}=
\sum\frac{\partial^\sigma f(dx)}{\partial dx^\sigma}\,\delta x^\sigma
=
\sum\frac{\partial^\sigma g(dy)}{\partial dy^\sigma}\,\delta y^\sigma,\\
\vdots
\end{gathered}
\tag{2}
\end{equation}

Каждая поляра через применение вариационного процесса \(\bdelta\) дает основание для необрывающейся последовательности общих инвариантов:
\begin{equation}
\begin{gathered}
\bdelta^\rho f(x,dx)=\bdelta^\rho g(y,dy),\\
\vdots\\
\bdelta^\rho f_{\delta^\sigma}
=\bdelta^\rho\sum\frac{\partial^\sigma f(dx)}{\partial dx^\sigma}\,
\delta x^\sigma
=\bdelta^\rho\sum\frac{\partial^\sigma g(dy)}{\partial dy^\sigma}\,
\delta y^\sigma,\\
\vdots\\[0.2em]
(\rho=0,1,2,\ldots).
\end{gathered}
\tag{3}
\end{equation}
Соотношения (2) и (3) в силу (1) позволяют считывать законы преобразования для производных от \(f(x,dx)\), но далее это не понадобится. Из инвариантов (3) теперь линейной комбинацией можно получить «нормальную форму \(\rho\)-й вариации», которая при \(\rho=1\) встречается у Лагранжа, а при \(\rho=2\) у Римана; она характеризуется тем, что смешанные дифференциалы порядка \(\rho+1\): \(d^\rho\delta\), \(d^{\rho-1}\delta^2,\ldots\), элиминированы. Для нее получаем:
\begin{equation}
\begin{gathered}
\Omega_1=\delta f(dx)-d f_\delta(dx),\\
\Omega_2=\delta^2 f-\delta d f_\delta+\frac12 d^2 f_{\delta^2},\\
\vdots\\
\Omega_\rho=\delta^\rho f-\delta^{\rho-1}d f_\delta
+\frac12\delta^{\rho-2}d^2 f_{\delta^2}
+\cdots+\frac{(-1)^\rho}{\rho!}\,d^\rho f_{\delta^\rho},\\
\vdots
\end{gathered}
\tag{4}
\end{equation}

Из инвариантов (4), обобщая подход Римана, можно перейти к инвариантам, содержащим только первые дифференциалы; такие инварианты будут называться «основными функциями». А именно, если в \(\Omega_1\) заменить дифференциалы дифференциальными отношениями, то \(\Omega_1=0\), тождественно по \(\delta x\), дает ровно лагранжевы уравнения вариационной задачи, принадлежащей
\[
f\!\left(\frac{dx}{dt}\right).
\]
Теперь я ставлю инвариантное условие, что направление \((d+\lambda\delta)\) удовлетворяет этим лагранжевым уравнениям:
\begin{equation}
\Omega_1(d+\lambda\delta)=
\bdelta f(dx+\lambda\delta x)
-(d+\lambda\delta)f_{\bdelta}(dx+\lambda\delta x)=0
\quad(\text{тождественно по }\delta x),
\tag{5}
\end{equation}
откуда подстановкой \(\bdelta=d+\lambda\delta\) для однородного \(f(dx)\) еще следует
\begin{equation}
(d+\lambda\delta)f(dx+\lambda\delta x)=0,
\tag{6}
\end{equation}
за исключением случая однородности первого порядка. В последнем случае (6) следует добавить как новое условие, поскольку тогда система уравнений (5) представляет только \((n-1)\) независимых условий. Если в (5), соответственно в (5) и (6), приравнять к нулю коэффициенты при нулевой, первой и второй степенях \(\lambda\), то получаются три инвариантные системы уравнений
\(\varphi=0\), \(\psi=0\), \(\chi=0\), тождественно по \(\delta x\), посредством которых \(d^2x,ddx,\delta^2x\) можно выразить через первые дифференциалы. Дальнейшие инвариантные системы уравнений
\begin{equation}
\begin{gathered}
d^\sigma\varphi=0,\qquad d^\sigma\psi=0,\qquad
d^\sigma\chi=0,\qquad d^{\sigma-1}\delta\chi=0,\\
d^{\sigma-2}\delta^2\chi=0,\ldots,\delta^\sigma\chi=0
\qquad(\sigma=0,1,2,\ldots)
\end{gathered}
\tag{7}
\end{equation}
тогда как раз достаточны, чтобы выразить все высшие дифференциалы через первые.\footnote{Только линейная форма \(\sum A_i\,dx_i\) требует особого рассмотрения, поскольку здесь система уравнений (5) не содержит вторых дифференциалов.} Следовательно, функции (4), соединенные с инвариантной системой уравнений (7), ведут к основным функциям \([\Omega_\rho]\), содержащим только первые дифференциалы; при этом \([\Omega_1]\) тождественно исчезает.

Так же из дифференциала \(\delta h(dx,\delta x)\) основной функции \(h\) посредством (7) снова можно получить основную функцию, которую следует назвать «ковариантной производной» \([h^{(1)}]\) от \(h\); вообще \([h^{(\nu)}]\) будет обозначать ковариантную производную от \([h^{(\nu-1)}]\). Эти два процесса таким образом ведут к дважды бесконечной последовательности основных функций:
\begin{equation}
\begin{array}{cccccc}
[\Omega_2], & [\Omega_2^{(1)}], & [\Omega_2^{(2)}], & \ldots,
& [\Omega_2^{(\nu)}], & \ldots\\
\vdots &&&& \vdots&\\
{}[\Omega_\rho], & [\Omega_\rho^{(1)}], & \ldots, & \ldots,
& [\Omega_\rho^{(\nu)}], & \ldots\\
\vdots & \vdots &&& \vdots&
\end{array}
\tag{8}
\end{equation}

Далее оказывается, что левые части уравнений, выражающих высшие дифференциалы через первые, когредиентны первым дифференциалам, то есть подчиняются тем же линейным преобразованиям. Поэтому, если вместо высших дифференциалов ввести эти левые части \(p,q,r,\ldots\) как аргументы инварианта, то, кроме производных, он зависит только от взаимно когредиентных величин. Указанное выше образование функций (8) просто сводится к введению этих новых величин; каждый инвариант тем самым переходит в сумму инвариантов, и из нее выбирается тот, который содержит именно \(dx,\delta x\), но не содержит \(p,q,r,\ldots\).

Ниже показывается, что при предположении однородности
\[
f(\varkappa\cdot dx)=\Phi(\varkappa)\cdot f(dx)
\]
основные функции (8) исчерпывают искомую полную систему.

\begin{center}
\textbf{II. Нормальные координаты, эквивалентность и теорема редукции.}
\end{center}

К нормальным координатам я прихожу интегрированием вариационной задачи, принадлежащей \(f\!\left(\frac{dx}{dt}\right)\):
\begin{equation}
\begin{gathered}
\Omega_1=\delta f\!\left(\frac{dx}{dt}\right)
-\frac{d}{dt}f_\delta\!\left(\frac{dx}{dt}\right)=0
\quad(\text{тождественно по }\delta x),\\
\frac{d}{dt}f\!\left(\frac{dx}{dt}\right)=0
\quad(\text{как следствие или дополнительное условие}).
\end{gathered}
\tag{9}
\end{equation}
Из-за предполагаемой однородности \(f(dx)\) эта система уравнений переходит в себя при подстановке параметра \(t=\varkappa\cdot\tau\); следовательно, если в силу (9) выразить \(\frac{d^\rho x_i}{dt^\rho}\) как функцию \(\varphi_\rho^{(i)}\!\left(\frac{dx}{dt}\right)\) первых производных, то \(\varphi_\rho^{(i)}\) становится однородной \(\rho\)-го порядка:
\begin{equation}
\varphi_\rho^{(i)}\!\left(\varkappa\cdot\frac{dx}{dt}\right)
=\varkappa^\rho\varphi_\rho^{(i)}\!\left(\frac{dx}{dt}\right).
\tag{10}
\end{equation}
Если теперь при интегрировании (9) для \(t=t_0\) задать начальные значения
\[
x=(x)_0,\qquad \frac{dx}{dt}=\left(\frac{dx}{dt}\right)_0
\]
и положить
\begin{equation}
u_i=(t-t_0)\left(\frac{dx_i}{dt}\right)_0,
\tag{11}
\end{equation}
где параметр \((t-t_0)\) в силу \(f\!\left(\frac{dx}{dt}\right)=\mathrm{const.}\) становится пропорциональным «длине экстремали», то разложение по степеням \((t-t_0)\) в силу (10) дает
\begin{equation}
x_i-(x_i)_0=u_i+\frac12\varphi_2^{(i)}(u)+\cdots+
\frac{1}{\rho!}\varphi_\rho^{(i)}(u)+\cdots .
\tag{12}
\end{equation}
Это разложение можно рассматривать как преобразование переменных \(x_i=x_i(u)\), где величины \(u\) как раз являются римановыми нормальными координатами. Значит, по (11) и (12) эти координаты преобразуют экстремали, проходящие через \((x)_0\), в прямые. Кроме того, произвольному преобразованию переменных \(x=x(y)\) по (11) соответствует линейное преобразование \(u=u(v)\) соответствующих нормальных координат, причем коэффициенты зависят только от произвольной системы значений \((x)_0\); следовательно, дифференциалы \(du,\delta u\) подчиняются тому же преобразованию, что и сами \(u\).

В силу (12) пусть \(f(x,dx)\) переходит в \(F(u,du)\), или, разложенная по степеням \(u\),
\begin{equation}
F(u,du)=F_0(u,du)+F_1(u,du)+\cdots+F_\rho(u,du)+\cdots
\tag{13}
\end{equation}
где \(F_\rho(u,du)\) обозначает однородную форму \(\rho\)-й размерности по \(u\). Из-за линейности преобразования \(u=u(v)\) из \(F(u,du)=G(v,dv)\) отдельные однородные составные части также переходят в соответствующие:
\begin{equation}
F_\rho(u,du)=G_\rho(v,dv).
\tag{14}
\end{equation}
Итак, если взять для \((x)_0\) какую-нибудь фиксированную постоянную систему значений, то (13) и (14) показывают, что эквивалентность \(f(dx)\) с \(g(dy)\) равносильна эквивалентности соответствующих систем функций \(F_\rho(u,du)\) относительно линейного преобразования. Кроме того, \(F_\rho(u,du)\), или соответствующие \(F_\rho(\delta u,du)\), представляют инварианты \(f(dx)\), взятые в точке \(x=(x)_0\), и именно они образуют полную систему основных функций для этой точки \(x=(x)_0\). А именно, имеем
\[
F_\rho(\delta u,du)=\frac1{\rho!}\sum
\frac{\partial F_\rho}{\partial \delta u^\rho}\,\delta u^\rho;
\qquad
\frac{\partial F_\rho}{\partial \delta u^\rho}
=\left(\frac{\partial F(du)}{\partial u^\rho}\right)_0,
\]
и эта последняя производная в силу (12) выражается через производные от \(f(dx)\), взятые при \(x=(x)_0\), совершенно так же, как соответственно образованная производная от \(G(dv)\) выражается через производные от \(g(dy)\). Поскольку
\[
\left(\frac{\partial^{\rho+\sigma}F(du)}
{\partial u^\rho\partial du^\sigma}\right)_0
=
\frac{\partial^{\rho+\sigma}F_\rho(\delta u,du)}
{\partial\delta u^\rho\partial du^\sigma}
\]
каждый инвариант \(f(dx)=F(du)\), взятый в точке \(x=(x)_0\), становится проективным инвариантом \(F_\rho(\delta u,du)\), как только высшие дифференциалы заменены на \(p,q,r,\ldots\), когредиентные к \(dx,\delta x\). Но \((x)_0\) можно рассматривать как параметр; каждому инварианту в точке \(x=(x)_0\) соответствует инвариант \(f(dx)\), если только снова заменить \((x)_0\) на \(x\). Следовательно, если \(F_\rho(\delta u,du)\) становятся инвариантами \(\Psi_\rho(\delta x,dx)\), взятыми при \(x=(x)_0\), ибо при \(x=(x)_0\) величины \(du,\delta u\) переходят в \(dx,\delta x\), то сами \(\Psi_\rho\) образуют основные функции полной системы. Остается выразить эту систему основных функций через явные инварианты \(f(dx)\), а именно через функции (8). Возможность этого показывает их порождение процессами дифференцирования. Для членов высшего порядка эти процессы переходят в простые процессы поляризации, и преобразование, аналогичное разложению в ряд Клебша--Гордана, вместе с индукционным выводом относительно отброшенных членов доказывает утверждение. Если речь идет об инвариантах одновременной системы, то к основным функциям (8) нужно еще добавить ковариантные производные остальных дифференциальных выражений, как показывает введение в эти выражения нормальных координат, принадлежащих \(f(dx)\).

Эквивалентность \(f(dx)\) с \(g(dy)\), которая была сведена к (14), следовательно, также равносильна эквивалентности \(\Psi_\rho\), или равным образом функций (8), относительно линейного преобразования дифференциалов, без необходимости в особых условиях интегрируемости, как следует из возможности сведения к (14). Тем самым теорема редукции доказана во всех частях. В частности, тождественное исчезновение функциональной последовательности \([\Omega_2],[\Omega_3],\ldots,[\Omega_\rho],\ldots\) необходимо и достаточно для того, чтобы \(f(dx)\) можно было преобразовать в выражение с постоянными коэффициентами.

Наконец, если вместо группы всех аналитических преобразований положить в основу подгруппу, то и группа соответствующих линейных преобразований \(u\) переходит в подгруппу проективной группы; инварианты \(f(dx)\) снова становятся инвариантами функций (8) относительно линейного преобразования, но теперь относительно этой подгруппы. Так, посредством однороднения, случай неоднородных функций \(f(dx)\) сводится к аффинной группе; полную систему здесь можно вывести из функций (8), образованных для одной дополнительной переменной.

Из теоремы редукции получается специальный случай квадратичных форм, более общо форм \(p\)-й размерности: здесь система функций обрывается на \([\Omega_2]\), более общо на \([\Omega_p]\) и ее ковариантных производных, поскольку все последующие основные функции становятся их проективными инвариантами. Этим объясняется выдающееся положение римановой формы кривизны \([\Omega_2]\) при квадратичных формах. Вопрос об эквивалентности квадратичных форм, соответственно форм \(p\)-й размерности, именно сводится к эквивалентности \([\Omega_2]\), соответственно \([\Omega_2]\ldots[\Omega_p]\), и их ковариантных производных относительно линейного преобразования; а тождественное исчезновение \([\Omega_2]\), соответственно \([\Omega_2]\ldots[\Omega_p]\), необходимо и достаточно для преобразования в формы с постоянными коэффициентами, чем для квадратичных форм высказан один из наиболее известных результатов.

Более подробное изложение должно вскоре появиться в Math. Annalen.

\clearpage

\end{document}
